\documentclass[11pt,a4paper]{article}

\usepackage{cmap}
\usepackage[T1]{fontenc}
\usepackage{lmodern}
\usepackage[utf8]{inputenc}
\usepackage[english]{babel}
\usepackage{fontawesome5}
\usepackage{amsmath,amssymb,mathtools}
\usepackage[a4paper,top=20mm,bottom=21mm,left=21mm,right=21mm,headheight=16pt]{geometry}
\usepackage{booktabs,longtable,tabularx,array,multicol}
\usepackage[table]{xcolor}
\usepackage{enumitem}
\usepackage[ruled,vlined]{algorithm2e}
\usepackage{fancyhdr}
\usepackage{hyperref}
\usepackage[expansion=false]{microtype}
\usepackage{titlesec}
\usepackage{pifont}
\usepackage{wasysym}
\usepackage{xparse}


\definecolor{siteblue}{HTML}{173B8F}
\definecolor{sitegold}{HTML}{F2C94C}
\definecolor{muted}{HTML}{64748B}
\definecolor{rulegray}{HTML}{B8C0CC}
\definecolor{softgray}{HTML}{F4F5F7}

\hypersetup{
  unicode=true,
  pdfencoding=auto,
  colorlinks=true,
  linkcolor=siteblue,
  urlcolor=siteblue,
  citecolor=siteblue,
  pdfauthor={Martin Minchev},
  pdftitle={Probability and Statistics: Alternative Exercises, 2026/2027}
}

\pagestyle{fancy}
\fancyhf{}
\lhead{2026/2027}
\rhead{\nouppercase{\leftmark}}
\cfoot{\thepage}
\renewcommand{\headrulewidth}{0.4pt}
\renewcommand{\sectionmark}[1]{\markboth{#1}{#1}}

\titleformat{\section}{\normalfont\large\bfseries}{\thesection.}{0.55em}{}
\titleformat{\subsection}{\normalfont\normalsize\bfseries}{}{0pt}{}
\titlespacing*{\section}{0pt}{1.2em}{0.55em}
\titlespacing*{\subsection}{0pt}{0.85em}{0.3em}
\newcommand{\sectionbreak}{\clearpage}

\setlist[itemize]{leftmargin=1.55em,itemsep=0.18em,topsep=0.25em}
\setlist[enumerate]{leftmargin=2.2em,itemsep=0.35em,topsep=0.3em}
\newenvironment{enumeratea}{\begin{enumerate}[label=(\alph*),itemsep=0.28em,topsep=0.25em]}{\end{enumerate}}
\setlength{\parindent}{0pt}
\setlength{\parskip}{0.32em}
\setlength{\emergencystretch}{3em}
\renewcommand{\arraystretch}{1.15}
\raggedbottom


\newcommand{\E}{\mathbb E}
\renewcommand{\P}{\mathbb P}
\newcommand{\Pbb}{\mathbb P}
\newcommand{\Ebb}{\mathbb E}
\newcommand{\Var}{\operatorname{Var}}
\newcommand{\Cov}{\operatorname{Cov}}
\newcommand{\Cor}{\operatorname{Cor}}
\newcommand{\Ber}{\operatorname{Ber}}
\newcommand{\Bin}{\operatorname{Bin}}
\newcommand{\Poi}{\operatorname{Poi}}
\newcommand{\Geom}{\operatorname{Geom}}
\newcommand{\Exp}{\operatorname{Exp}}
\newcommand{\Unif}{\operatorname{U}}
\newcommand{\ind}{\mathbf 1}
\newcommand{\iid}{\mathrel{\mathrm{iid}}}

\newcounter{problem}[section]
\renewcommand{\theproblem}{\thesection.\arabic{problem}}
\newcommand{\problemstart}[2]{%
  \refstepcounter{problem}%
  \par\medskip\noindent
  #1\if\relax\detokenize{#1}\relax\else\,\fi
  \textbf{Problem \theproblem.}%
  \if\relax\detokenize{#2}\relax\else\ \emph{(#2)}\fi
  \enspace\ignorespaces
}
\NewDocumentEnvironment{problem}{O{}}{\problemstart{}{#1}}{\par}
\NewDocumentEnvironment{problembook}{O{}}{\problemstart{\faBook}{#1}}{\par}
\NewDocumentEnvironment{problemstar}{O{}}{\problemstart{\faStar}{#1}}{\par}
\NewDocumentEnvironment{probleminteresting}{O{}}{\problemstart{\faSmile[regular]}{#1}}{\par}
\NewDocumentEnvironment{problempencil}{O{}}{\problemstart{\faPencil*}{#1}}{\par}
\NewDocumentEnvironment{problemstarpencil}{O{}}{\problemstart{\faStar\,\faPencil*}{#1}}{\par}
\NewDocumentEnvironment{researchtask}{O{}}{\problemstart{\faRobot}{#1}}{\par}

\newcommand{\aitag}[1]{{\footnotesize\sffamily\color{muted}AI--#1}}
\newcommand{\projectcheck}[1]{%
  \par\smallskip\noindent\colorbox{softgray}{%
    \parbox{\dimexpr\linewidth-2\fboxsep}{\small\textbf{Project checkpoint.} #1}}\par\smallskip
}
\newcommand{\leancheck}[1]{%
  \par\smallskip\begingroup\sloppy\noindent{\small\textbf{Lean micro-task.} #1}\par\endgroup\smallskip
}

\newcommand{\chatgptmodel}{\textit{[model and version]}}
\newcommand{\claudemodel}{\textit{[model and version]}}

\begin{document}
\setcounter{tocdepth}{1}
\tableofcontents
\noindent{\footnotesize English translation prepared with ChatGPT (OpenAI, GPT-6),
26 October 2026, for translation of the Bulgarian text into English and LaTeX formatting.
The Bulgarian originals remain the reference; translation may contain errors.}\par\smallskip

\subsection*{Notes}
Each topic contains \emph{core problems}, covering the minimum required for that topic;
\emph{additional problems}, providing variety and extending the core material; and
\emph{challenges}, related to the topic but going beyond the requirements of standard
introductory courses in probability and statistics.
The following symbols are used:
\begin{itemize}
  \item[\faBook] a problem introducing fundamental objects of the course;
  \item[\faStar] a challenging problem;
  \item[\faSmile] an interesting example;
  \item[\faPencil*] a problem from an FMI test;
  \item[\faRobot] an independent research task. All resources are permitted, including
  literature, software and AI. You are expected to check the results, understand them,
  and be able to explain them.
\end{itemize}
{\small
\textbf{Note on the use of AI.}
The latest publicly available versions of ChatGPT and Claude (at the time of writing,
sol 6.1 and opus 5.5, respectively) were used for formatting, editing and clarifying
individual formulations. The compiler is responsible for the final version and its
mathematical correctness.
}
\section{Combinatorics}
\subsection*{Core problems}
\begin{problembook}
For $n,k\ge 1$, recall:
\begin{enumeratea}
  \item how many solutions does $x_1+\cdots+x_n=k$ have when $x_1,\dots,x_n$
  are positive integers? What if they are nonnegative integers?
  \item state the inclusion--exclusion principle.
\end{enumeratea}
\end{problembook}
\begin{problembook}
Let $M$ be a set with $n$ elements. In how many ways can we choose:
\begin{enumeratea}
  \item $k$ distinct elements of $M$;
  \item an ordered $k$-tuple of distinct elements of $M$;
  \item an ordered $k$-tuple of elements of $M$, with repetitions allowed?
\end{enumeratea}
\end{problembook}
\begin{problem}
In how many ways can all $k$ distinguishable particles be distributed among $n$
distinct cells if:
\begin{enumeratea}
  \item each cell contains at most one particle;
  \item cells may contain any number of particles;
  \item no cell is empty?
\end{enumeratea}
Answer the same questions for indistinguishable particles. Also state when the answer is zero.
\end{problem}
\begin{problem}
Five distinguishable balls are distributed among three distinct boxes $A,B,C$.
Find the number of distributions in which:
\begin{enumeratea}
  \item $A$ is empty;
  \item only $A$ is empty;
  \item exactly one box is empty;
  \item at least one box is empty;
  \item no box is empty.
\end{enumeratea}
\end{problem}
\begin{problemstar}
Consider the following questions.
\begin{enumeratea}
  \item In how many ways can we go from $(0,0)$ to $(n,n)$ on the square lattice
  using unit steps to the right and upwards? What if the path never goes above the diagonal?
  \item How many correctly balanced parenthesis expressions have $n$ opening and $n$
  closing parentheses? Construct an explicit bijection between them and the paths
  in the second part of the preceding question.
  \item Find simple asymptotic equivalents for the total number of paths and for the
  number that never go above the diagonal, as $n\to\infty$.
\end{enumeratea}
\end{problemstar}
\subsection*{Additional problems}
\begin{problem}
How many four-digit numbers can be formed from the digits $1,2,3,4,5$ if:
\begin{enumeratea}
  \item repetition is not allowed;
  \item repetition is allowed;
  \item repetition is not allowed and the number is odd?
\end{enumeratea}
\end{problem}
\begin{problem}
In how many ways can a four-person delegation be selected from 12 candidates,
including distinct candidates $A,B,C,D$, if:
\begin{enumeratea}
  \item there are no restrictions;
  \item $A$ and $B$ must not both participate;
  \item either both $C$ and $D$ participate, or neither does?
\end{enumeratea}
\end{problem}
\begin{problem}
Prove combinatorially, and then algebraically, that
\[
\sum_{k=0}^n\binom nk=2^n
\qquad\text{and}\qquad
\binom{2n}{n}=\sum_{k=0}^n\binom nk^2.
\]
\end{problem}
\begin{problem}
Let $p(n)$ be the number of partitions of $n$ into distinct positive integers,
and $q(n)$ the number of partitions into odd positive integers, where repetitions
are allowed and the order of the summands does not matter. For example, for $n=6$:
\[
p(6)=4:\ 6,\ 5+1,\ 4+2,\ 3+2+1; \qquad
q(6)=4:\ 5+1,\ 3+3,\ 3+1+1+1,\ 1+1+1+1+1+1.
\]
\begin{center}
\begin{tabular}{c|cccccccc}
  $n$    & 1 & 2 & 3 & 4 & 5 & 6 & 7 & 8 \\ \hline
  $p(n)$ & 1 & 1 & 2 & 2 & 3 & 4 & 5 & 6 \\
  $q(n)$ & 1 & 1 & 2 & 2 & 3 & 4 & 5 & 6
\end{tabular}
\end{center}
By comparing generating functions, prove that $p(n)=q(n)$ for every $n$.
\end{problem}
\subsection*{Challenges and independent research}
\begin{problemstar}[Derangements]
Let $D_n$ be the number of permutations $\sigma\in S_n$ such that $\sigma(i)\ne i$ for every $i$.
\begin{enumeratea}
  \item Prove that
  \[ D_n=n!\sum_{j=0}^{n}\frac{(-1)^j}{j!}. \]
  \item Prove combinatorially that, for $n\ge 2$,
  \[ D_n=(n-1)(D_{n-1}+D_{n-2}), \]
  where $D_0=1$ and $D_1=0$.
  \item Prove that, for $n\ge 1$, $D_n$ is the nearest integer to $n!/e$,
  without using a numerical approximation to $e$.
\end{enumeratea}
\end{problemstar}
\begin{problemstar}[Necklaces and Burnside's lemma]
Let $X=\{0,1\}^n$ be the set of binary words of length $n$, and let
$G=\{r^0,r^1,\dots,r^{n-1}\}\cong\mathbb Z_n$ be the group of cyclic rotations,
$r(x_1x_2\dots x_n)=x_nx_1\dots x_{n-1}$. A \emph{necklace} is an orbit of $G$
in $X$: two words are identified if one is obtained from the other by a cyclic
rotation. For example, $0011,\,0110,\,1100,\,1001$ form one necklace, and
$0101,\,1010$ form another. Rotations preserve the number of ones, so $G$ also acts on
\[ X_k=\{x\in X:\ x\text{ has exactly }k\text{ ones}\}. \]
Write $N(n,k)$ for the number of orbits in $X_k$, that is, the number of necklaces
with $n$ beads and exactly $k$ ones.
For $g\in G$, let
\[ \operatorname{Fix}(g)=\{x\in X_k:\ g(x)=x\}, \]
and let $X_k/G$ be the set of orbits. Use Burnside's (Cauchy--Frobenius) lemma
\[ \#(X_k/G)=\frac1{|G|}\sum_{g\in G}|\operatorname{Fix}(g)|. \]
Explain why a rotation $g$ whose cycles on the positions $\{1,\dots,n\}$ have
length $d$ can fix a word with $k$ ones only if $d\mid k$. Prove that in this case
\[ |\operatorname{Fix}(g)|=\binom{n/d}{k/d}, \]
whereas if $d\nmid k$, then $\operatorname{Fix}(g)=\emptyset$. Also prove that
for every $d\mid n$ there are exactly $\varphi(d)$ rotations whose cycles have length $d$.
Deduce that, for $1\le k\le n-1$,
\[ N(n,k)=\frac1n\sum_{d\mid\gcd(n,k)}\varphi(d)\binom{n/d}{k/d}, \]
where $\varphi$ is Euler's totient function. As a check, find $N(12,4)$.
\end{problemstar}
\begin{researchtask}[Random mappings and cyclic vertices]
Let $f:[n]\to[n]$ be chosen uniformly from all $n^n$ mappings. Let $G_f$ be the
directed graph with edges $i\to f(i)$, and $K_n$ the number of vertices lying on
a cycle in $G_f$. You may use AI to propose solutions, acting as their \emph{reviewer}:
verification, rather than generation, is assessed.
\begin{enumeratea}
  \item Use software to find the distribution $\Pbb(K_n=k)$ for $n\le 7$.
  The code may be AI-generated, but run it and check that
  $\sum_k\#\{K_n=k\}=n^n$; also check the case $n=3$ by hand.
  \item Find and prove a formula for $\#\{f:K_n=k\}$. You may ask AI for a proposal.
  (i) Compare it with the table from part 1. (ii) Write the proof in your own words,
  separating the choice of cyclic vertices, the permutation on them, and the directed
  forest rooted at them. (iii) Give your own proof of the formula $k\,n^{n-k-1}$
  for the number of forests with $k$ prescribed roots. If you use AI, check whether
  it has actually proved the formula, and give a proof by a different method.
  (iv) Find and numerically check the number of mappings with $K_n=k$ and exactly $c$ cycles.
  \item Derive $\Ebb K_n$ in two ways: from the distribution and by linearity,
  computing $\Pbb(\text{vertex }1\text{ is cyclic})$. Prove the identity between
  the two expressions. Use the exact formula to calculate $\Ebb K_n$ for
  $n=10,10^2,10^3,10^4,10^5$ and investigate numerically whether
  \[ \frac{\Ebb K_n}{\sqrt n}\longrightarrow\sqrt{\frac{\pi}{2}}. \]
  Explain the origin of the constant $\sqrt{\pi/2}$. Optionally investigate the finer asymptotic claim
  \[ \Ebb K_n-\sqrt{\frac{\pi n}{2}}\longrightarrow-\frac13. \]
  Let $T_n$ be the number of distinct vertices in the trajectory $1,f(1),f^2(1),\dots$
  before its first repetition. Do $K_n$ and $T_n$ have the same distribution,
  or only the same expectation?
\end{enumeratea}
\end{researchtask}
\noindent\textbf{Lean.} Let $A,B\subseteq\Omega$, with indicators taking values in $\mathbb N$.
Define the indicator by
\[
\ind_A(x)=\begin{cases}1,&x\in A,\\0,&x\notin A,\end{cases}
\]
and prove in Lean 4 with Mathlib that, for every $x\in\Omega$,
\[ \ind_{A\cup B}(x)+\ind_{A\cap B}(x)=\ind_A(x)+\ind_B(x). \]

\section{Conditional probability, independence and Bayes' rule}
\subsection*{Core problems}
\begin{problembook}
Two standard dice are rolled independently.
\begin{enumeratea}
  \item What is the probability that the sum is less than 8, given that it is odd?
  Are the events ``the sum is less than 8'' and ``the sum is odd'' independent?
  \item Are ``the sum is 7'' and ``the first die shows 3'' independent?
  What about ``the sum is 8'' and ``the first die shows 3''?
\end{enumeratea}
\end{problembook}
\begin{problem}
Let $A$ and $B$ be events with $\Pbb(A)>0$ and $0<\Pbb(B)<1$.
\begin{enumeratea}
  \item Prove that if $A$ and $B$ are independent, then so are $A$ and $\overline{B}$,
  and also $\overline{A}$ and $\overline{B}$.
  \item Prove that the following four conditions are equivalent:
  \[
  \Pbb(A\cap B)>\Pbb(A)\Pbb(B),\quad
  \Pbb(A\mid B)>\Pbb(A),\quad
  \Pbb(B\mid A)>\Pbb(B),\quad
  \Pbb(A\mid B)>\Pbb(A\mid \overline{B}).
  \]
  \item If $A$ is independent of $B$ and $A$ is independent of $C$, must $A$ be
  independent of $B\cap C$? What about $B\cup C$?
\end{enumeratea}
\end{problem}
\begin{problembook}
Three cards are drawn successively without replacement from a deck of 52 cards.
Find the probability that:
\begin{enumeratea}
  \item all three are aces;
  \item the third card is an ace;
  \item the second card is an ace, given that the first is not.
\end{enumeratea}
\end{problembook}
\begin{problembook}
A box contains three coins: two fair coins and one with heads on both sides.
We choose a coin uniformly and toss it. Conditional on the chosen coin, successive tosses are independent.
\begin{enumeratea}
  \item The outcome is heads. What is the probability that the coin is the false one?
  \item We toss the same coin again and get heads again. Find the new probability
  in two ways: by updating the preceding answer, treating it as a new prior,
  and directly from the two tosses.
  \item What is the probability of heads on the third toss, given that the first two were heads?
\end{enumeratea}
\end{problembook}
\begin{problempencil}[Test 1, SE 2023]
An urn contains 4 yellow, 5 green and 2 blue pebbles. We select a pebble uniformly.
If it is yellow, we roll one standard die; if green, two; if blue, three.
The rolls are independent. Given that the sum is 3, what is the probability that the pebble was green?
\end{problempencil}
\begin{problembook}
A test gives a positive result for a person with a disease with probability $99\%$
(sensitivity), and a negative result for a healthy person with probability $99\%$
(specificity). The disease affects $0.5\%$ of the population. We choose a person uniformly and test them.
\begin{enumeratea}
  \item What is the probability that a person with a positive result actually has the disease?
  \item At what disease prevalence is this probability exactly $1/2$?
  \item The person is tested again and the result is again positive. If the two tests
  are conditionally independent given the person's health status, and have the same
  sensitivity and specificity, what is the probability that the person has the disease?
  When is this independence assumption questionable?
\end{enumeratea}
\end{problembook}
\subsection*{Additional problems}
\begin{probleminteresting}[Kahneman, Thinking, Fast and Slow]
An accident occurs at night in a city with a blue and a green taxi company.
$85\%$ of the taxis are green and $15\%$ are blue. Assume that a taxi's probability
of being involved in an accident does not depend on its company.
A witness states that the taxi was blue. A subsequent assessment shows that the
witness identifies the colour correctly with probability $80\%$, both for blue
and for green taxis. What is the probability that the taxi really was blue?
\end{probleminteresting}
\begin{problem}
Three tokens have sides coloured white--white, black--black and white--black,
respectively. We choose a token uniformly and place it with a randomly selected
side facing up, the two sides being equally likely. Given that the upper side is
white, what is the probability that the hidden side is also white?
\end{problem}
\begin{probleminteresting}[Monty Hall]
A car is behind one of three doors, chosen uniformly. You select a door without
knowing where the car is. The host knows its location and always opens a different
empty door, then offers you the chance to switch to the remaining closed door.
\begin{enumeratea}
  \item Should you switch? Does the probability of winning under the ``always switch''
  strategy depend on how the host chooses a door when there are two possibilities?
  \item You chose door 1 and the host opened door 3. When the host can choose between
  doors 2 and 3, they open door 3 with probability $p\in[0,1]$.
  Find the conditional probability that the car is behind door 2. For which $p$ is switching worthwhile?
\end{enumeratea}
\end{probleminteresting}
\begin{probleminteresting}[Birthday paradox]
Assume that different people's birthdays are independent and uniformly distributed
over 365 days. What is the smallest number of people for which the probability
that at least two share a birthday is greater than $1/2$? Use the multiplication
rule. Discuss the assumptions of the model.
\end{probleminteresting}
\begin{problem}
Two fair coins are tossed independently. Let $A$ be ``the first coin shows heads'',
$B$ be ``the second coin shows heads'', and $C$ be ``the coins show the same outcome''.
Prove that $A,B,C$ are pairwise independent but not mutually independent.
Are $A$ and $C$ independent conditional on $B$?
\end{problem}
\begin{probleminteresting}[Simpson's paradox]
The following table contains real data on two treatments for kidney stones
(Charig et al., \emph{BMJ} 292, 1986; Julious \& Mullee, \emph{BMJ} 309, 1994):
\begin{center}\small
\begin{tabular}{@{}lrr@{}}
\toprule
 & Treatment A & Treatment B \\
\midrule
Small stones & $93\%\ (81/87)$ & $87\%\ (234/270)$ \\
Large stones & $73\%\ (192/263)$ & $69\%\ (55/80)$ \\
Overall & $78\%\ (273/350)$ & $83\%\ (289/350)$ \\
\bottomrule
\end{tabular}\end{center}
Explain why the comparison between the treatments reverses when the groups are combined.
Can these data alone establish which treatment is better?
\end{probleminteresting}
\begin{probleminteresting}[Reservoir sampling]
Let $A$ be an array of $n\ge3$ distinct elements, indexed starting at 1.
The following algorithm is executed, with independent random choices in different iterations:
\begin{algorithm}[H]
\caption{Reservoir sampling for one element}
$sample \gets A[1]$\;
\For{$i \gets 2$ \KwTo $n$}{
  $j \gets$ a random integer, uniformly distributed on $\{1,\dots,i\}$\;
  \If{$j = 1$}{ $sample \gets A[i]$\; }
}
\Return{$sample$}
\end{algorithm}
What is the probability that the final value of $sample$ is $A[1]$? $A[3]$? $A[n]$?
Prove that the algorithm selects an element of the array uniformly.
\end{probleminteresting}
\begin{problem}
There are $n\ge1$ urns, each containing $m$ white and $k$ black balls, where $m+k>0$.
We transfer a uniformly chosen ball from the first urn to the second; after adding it,
we choose a ball uniformly from the current contents of the second urn and transfer
it to the third, and so on. Find the probability that a ball drawn uniformly from
the final urn at the end is white.
\end{problem}
\subsection*{Challenges and independent research}
\begin{problemstar}[A hidden common cause]
Let $E$ be an event with $\Pbb(E)=q$, where $0<q<1$, and let $A$ and $B$ be
conditionally independent given the value of $\ind_E$. Let
\[
\Pbb(A\mid \ind_E=e)=a_e
\quad\text{and}\quad
\Pbb(B\mid \ind_E=e)=b_e,
\qquad\text{for }e\in\{0,1\}.
\]
\begin{enumeratea}
  \item Prove that
  \[ \Cov(\ind_A,\ind_B)=q(1-q)(a_1-a_0)(b_1-b_0). \]
  When is this covariance zero? Explain why, for indicators, this is equivalent
  to independence of $A$ and $B$.
  \item Assuming $\Pbb(A\cap B)>0$, find $\Pbb(E\mid A\cap B)$.
  \item Let
  \[ q=0.10,\qquad a_0=0.02,\quad a_1=0.20,\qquad b_0=0.03,\quad b_1=0.25. \]
  Calculate $\Pbb(A)$, $\Pbb(B)$ and $\Cov(\ind_A,\ind_B)$. Find and compare
  \[ \Pbb(E),\qquad\Pbb(E\mid A),\qquad\Pbb(E\mid B),\qquad\Pbb(E\mid A\cap B). \]
  Check that
  \[ \Pbb(A\cap B)=0.00554\qquad\text{and}\qquad\Pbb(E\mid A\cap B)=\frac{250}{277}. \]
  Interpret the results by viewing $E$ as a hidden common cause and $A,B$ as two observed clues to it.
\end{enumeratea}
\end{problemstar}
\begin{problemstar}[An urn of unknown composition]
For a fixed $n\ge2$, choose $K$ uniformly from $\{0,1,\dots,n\}$. Then use an
urn containing $K$ red and $n-K$ blue balls. Make two draws from it.
Consider separately sampling with and without replacement. Let $R_i$ be the event
``the $i$th ball drawn is red'', for $i\in\{1,2\}$.
\begin{enumeratea}
  \item Given that the first ball is red, prove that in both models
  \[ \Pbb(K=k\mid R_1)=\frac{2k}{n(n+1)},\qquad\text{for }k\in\{0,\dots,n\}. \]
  \item Find $\Pbb(R_2\mid R_1)$ in both models and prove that the results are, respectively,
  \[ \frac{2n+1}{3n}\qquad\text{and}\qquad\frac23. \]
  \item Explain probabilistically why the results differ for finite $n$ and why
  the difference disappears as $n\to\infty$.
\end{enumeratea}
\end{problemstar}
\begin{researchtask}[When does a repeated test fail to provide independent evidence?]
Let $D\sim\Ber(0.01)$ indicate the presence of a disease, and let $B\sim\Ber(q)$
indicate a problematic laboratory batch, with $D$ and $B$ independent. All tests
on a given person use the same batch. Conditional on $(D,B)$, the results
$T_1,T_2,\dots$ are independent and identically distributed, with $T_i=1$ meaning a positive result.
Initially let $q=0.10$, with positive-result probabilities given by
\[
\begin{array}{c|cc}
 & B=0 & B=1\\ \hline
\Pbb(T_i=1\mid D=1,B) & 0.95 & 0.80\\
\Pbb(T_i=1\mid D=0,B) & 0.05 & 0.30
\end{array}
\]
\begin{enumeratea}
  \item Derive an exact formula for
  \[ p_k=\Pbb(D=1\mid T_1=\cdots=T_k=1),\qquad\text{for }k\ge1. \]
  \item Compare it with the naive model that first averages sensitivity and the
  false-positive probability over $B$, then assumes that test results are
  conditionally independent given $D$. Explain why this assumption can change the answer.
  \item Make a table for $k=1,\dots,5$. Check that
  \[ p_1=\frac{17}{152},\qquad p_2=\frac{701}{1592}\approx0.4403, \]
  whereas the naive model gives approximately $0.6109$ for $k=2$.
  \item Consider the parameter grid
  \[
  q\in\{0,0.05,0.10,0.20\},\qquad
  \Pbb(T_i=1\mid D=0,B=1)\in\{0.10,0.20,0.30,0.40\},
  \]
  keeping all other parameters fixed. Find the pairs for which the naive estimate
  of $p_2$ exceeds the exact value by at least $0.10$.
  \item Does the naive model always overestimate $p_2$ if all four probabilities
  in the table may vary? Give a proof or a counterexample and explain the result.
  \item Does the naive model always overestimate $p_2$ if all four probabilities
  in the table may vary? Find a necessary and sufficient condition for overestimation,
  give a counterexample, and interpret the condition.
\end{enumeratea}
\end{researchtask}
\noindent\textbf{Lean.} Let $A,B\subseteq\Omega$, with indicators taking values
in $\mathbb N$. Use the definition from the first topic and prove in Lean 4 with
Mathlib that, for every $x\in\Omega$,
\[ \ind_{A\cap B}(x)=\ind_A(x)\ind_B(x). \]

\section{Random variables and discrete distributions}
\subsection*{Core problems}
\begin{problembook}
A player stakes 5 leva and rolls two independent standard dice. Two sixes pay
100 leva; exactly one six pays 5 leva; no sixes pay nothing. Let $X$ be the player's
net gain. Find the distribution and expectation of $X$. Is the game fair?
\end{problembook}
\begin{problembook}
A person tosses a fair coin, taking a step forwards for heads and backwards for
tails. Toss outcomes are independent. What is the probability that after 10 tosses the person is:
\begin{enumeratea}
  \item at the starting point;
  \item at distance 2 from it;
  \item five steps ahead of it?
\end{enumeratea}
\end{problembook}
\begin{problembook}
Let $Z_1$ and $Z_2$ be the outcomes of two independent dice rolls, and let $X=Z_1+Z_2$.
\begin{enumeratea}
  \item The dice are standard. Find the distribution, expectation and variance of $X$.
  \item Let
  \[
  \Pbb(Z_i=1)=\Pbb(Z_i=6)=\frac14
  \quad\text{and}\quad\Pbb(Z_i=j)=\frac18,
  \qquad\text{for }i\in\{1,2\},\ j\in\{2,3,4,5\}.
  \]
  Find the distribution, expectation and variance of $X$, showing two ways of finding the expectation.
  \item For standard dice, let $Y=|X-7|$. Find the distribution of $Y$ and $\Ebb Y$.
  Compare $\Ebb Y$ with $|\Ebb X-7|$, and $\Ebb Y^2$ with $\Var X$.
\end{enumeratea}
\end{problembook}
\begin{problem}
A box contains 7 light bulbs, 3 of them defective. Four distinct bulbs are chosen
uniformly. Find the distribution and expectation of the number of defective bulbs
chosen. Also find the expectation using indicators, without summing over the whole
distribution. Find the variance using indicators as well. Are these indicators independent?
\end{problem}
\begin{problembook}
Two independent standard dice are rolled five times, with outcomes in different
rolls also independent. What is the probability that exactly 2 rolls have sum 6?
Find the expectation and variance of the number of such rolls.
\end{problembook}
\begin{problempencil}[Memorylessness; Test 1, SE 2023]
Let $X\sim\Geom(p)$, meaning $\Pbb(X=k)=p(1-p)^k$ for $k\in\{0,1,2,\dots\}$,
where $p\in(0,1]$.
\begin{enumeratea}
  \item Prove that for all integers $m,n\ge0$ with $\Pbb(X\ge n)>0$,
  \[ \Pbb(X\ge m+n\mid X\ge n)=\Pbb(X\ge m). \]
  \item Conversely, let $Y$ take values in $\{0,1,2,\dots\}$ and satisfy the above
  equality for every pair of integers $m,n\ge0$ with $\Pbb(Y\ge n)>0$.
  Prove that $Y\sim\Geom(p)$ for some $p\in(0,1]$.
\end{enumeratea}
\end{problempencil}
\begin{problembook}
Assume that the number of minor earthquakes in a region over $t$ months has a
Poisson distribution with parameter $2t$. What is the probability of fewer than
four earthquakes in three months?
\end{problembook}
\begin{probleminteresting}[St. Petersburg paradox]
A casino charges an entry fee of $A$ leva. We toss a fair coin until the first
heads, with independent outcomes. If heads first occurs on toss $n$, the prize
is $2^n$ leva. Find the expected prize. For which $A$ would you personally play?
What changes if the prize is $\min\{2^n,2^{30}\}$ leva? How does this result
compare with your answer about the values of $A$ for which you would play?
\end{probleminteresting}
\subsection*{Additional problems}
\begin{problem}
An urn contains 5 white and 3 black balls. We draw until a white ball appears.
Let $X$ be the number of black balls drawn. Find its distribution, expectation and variance for:
\begin{enumeratea}
  \item sampling without replacement;
  \item sampling with replacement.
\end{enumeratea}
\end{problem}
\begin{problem}
$80\%$ of home printers work without additional configuration, with independent
outcomes for different printers and weeks. A company sells 10 printers each week.
Find the probability that at least 9 work immediately; the probability that this
happens in five consecutive weeks; and the probability that the first week in
which fewer than 9 work immediately is week $k$. Calculate the last probability
for $k=3$ and $k=21$, and find the expected number of weeks until the first such week.
\end{problem}
\begin{problempencil}[Test 1, SE 2023]
\begin{enumeratea}
  \item Let $X_1,\dots,X_n$ be independent and identically distributed, with
  $F(m)=\Pbb(X_1\le m)$, and let $X_{\min}=\min\{X_1,\dots,X_n\}$ and
  $X_{\max}=\max\{X_1,\dots,X_n\}$. Express $\Pbb(X_{\min}\le m)$ and
  $\Pbb(X_{\max}\le m)$ in terms of $F(m)$ and $n$.
  \item Five standard dice are rolled. In each round, we roll all remaining dice
  simultaneously and remove those showing six. All rolls are independent. Let $K$
  be the number of rounds until all dice have been removed. What is the distribution
  of the number of rolls of one die? Find $\Pbb(K\le n)$ and $\Pbb(K=n)$ for every integer $n\ge1$.
\end{enumeratea}
\end{problempencil}
\begin{problem}
Let $X$ take values in $\{0,1,2,\dots\}$ and have finite expectation.
\begin{enumeratea}
  \item Prove that
  \[ \Ebb X=\sum_{k=1}^\infty\Pbb(X\ge k)=\sum_{n=0}^\infty\Pbb(X>n). \]
  \item Check the formula for $X\sim\Geom(p)$.
  \item Find $\Ebb K$ for the number of rounds $K$ in the preceding problem and
  compare it with the expected number of rolls of one die.
\end{enumeratea}
\end{problem}
\begin{problem}
Independent Bernoulli trials have success probability $p$, where $0<p<1$.
Find the probability that the $r$th success occurs on exactly trial $k+r$, for
integers $r\ge1$ and $k\ge0$. Let $X_i$ indicate success on trial $i$. Express
the event using $X_i$ and the partial sums $S_n=\sum_{i=1}^nX_i$, with $S_0=0$.
\end{problem}
\begin{problem}
A shot hits its target with probability $0.001$. There are 5000 independent shots,
and at least two hits are required. Find the exact probability and the Poisson
approximation. How large is the numerical difference?
\end{problem}
\begin{problem}
A submarine fires $n$ times; shots are independent and each torpedo hits with
probability $p$. On a hit, one of the $m$ compartments is selected uniformly,
independently of other hits. Find the probability that at least two distinct compartments are flooded.
\end{problem}
\begin{problem}[Banach's matchboxes]
A smoker carries two boxes containing $n$ matches each. Each time, the smoker
independently chooses one box with probability $1/2$ and takes a match.
On first reaching for an empty box, what is the probability that the other box
contains exactly $k\in\{0,\dots,n\}$ matches?
\end{problem}
\subsection*{Challenges and independent research}
\begin{problemstar}
Let $\pi$ be a uniformly random permutation of $\{1,\dots,n\}$, and let
\[ I_n=\#\{(i,j):1\le i<j\le n,\ \pi(i)>\pi(j)\}. \]
Prove that
\[ \Ebb I_n=\frac{n(n-1)}4\qquad\text{and}\qquad\Var(I_n)=\frac{n(n-1)(2n+5)}{72}. \]
\end{problemstar}
\begin{problemstar}[Coupon collector]
Let $X_1,X_2,\dots$ be independent and identically distributed, uniform on
$\{1,\dots,m\}$, where $m\ge2$, and let
\[ T=\inf\{n\ge1:\{X_1,\dots,X_n\}=\{1,\dots,m\}\}. \]
\begin{enumeratea}
  \item Find $\Pbb(T\le n)$ for $n\ge1$.
  \item Prove that
  \[ \Ebb T=mH_m,\qquad\Var(T)=m^2H_m^{(2)}-mH_m, \]
  where $H_m=\sum_{j=1}^m\frac1j$ and $H_m^{(2)}=\sum_{j=1}^m\frac1{j^2}$.
\end{enumeratea}
Where do you use independence of the waiting times, and where is linearity of expectation alone sufficient?
\end{problemstar}
\begin{researchtask}[Generalised coupon collector]
Let $X_1,X_2,\dots$ be independent and identically distributed random variables
taking values in $[m]=\{1,\dots,m\}$, where $m\ge2$, with
\[
\Pbb(X_n=i)=p_i>0,\qquad\text{for }i\in[m],
\qquad\text{and }\sum_{i=1}^m p_i=1.
\]
Let $T$ be the number of draws until all $m$ values have appeared.
\begin{enumeratea}
  \item Prove that
  \[
  \Ebb T=\sum_{\emptyset\ne A\subseteq[m]}\frac{(-1)^{|A|+1}}{\sum_{i\in A}p_i}
  \qquad\text{and}\qquad
  \Ebb T=\int_0^\infty\left(1-\prod_{i=1}^m\bigl(1-e^{-p_it}\bigr)\right)\,dt.
  \]
  Explain why the first formula is unsuitable for direct numerical calculation
  when $m=30$. Suggest a numerical method for the integral and a way to check its accuracy.
  \item Compare values from the integral formula with simulations for $m=10,30,100$
  in the following models:
  \begin{enumerate}[label=(\roman*)]
    \item uniform probabilities: $p_i=1/m$ for $i\in[m]$;
    \item probabilities
    \[ p_i=\frac{i^{-\alpha}}{\sum_{j=1}^m j^{-\alpha}},\qquad\text{for }i\in[m], \]
    considering both $\alpha=1$ and $\alpha=2$;
    \item one rare coupon:
    \[ p_1=m^{-2}\quad\text{and}\quad p_i=\frac{1-m^{-2}}{m-1},\qquad\text{for }i\in\{2,\dots,m\}. \]
  \end{enumerate}
  State the number of simulations and discuss the accuracy of the estimated mean.
  Explain which mechanism determines the scale of $\Ebb T$ in each model.
\end{enumeratea}
You may use AI for the proof and numerical method, provided you check the results
and explain how you checked them.
\end{researchtask}
\noindent\textbf{Lean.} Let $S$ be a finite nonempty set, $X,Y:S\to\mathbb R$,
and $a,b\in\mathbb R$. On the finite uniform probability space on $S$, define
\[ \Ebb X=\frac1{|S|}\sum_{s\in S}X(s). \]
Formalise in Lean 4 with Mathlib and prove linearity:
\[ \Ebb(aX+bY)=a\,\Ebb X+b\,\Ebb Y. \]
First prove, by induction on $S$ and without using \texttt{Finset.sum\_add\_distrib}, that
\[ \sum_{s\in S}\bigl(X(s)+Y(s)\bigr)=\sum_{s\in S}X(s)+\sum_{s\in S}Y(s); \]
\texttt{Finset.mul\_sum} and \texttt{Finset.sum\_div} are useful for the rest.
Why is nonemptiness of $S$ needed to define $\Ebb$, but not for the identity for sums?

\section{Dependence, sums and generating functions}
\subsection*{Part I. Joint distributions and conditional expectation}
\begin{problembook}
We independently roll two red and one blue standard dice. Let $X$ be the number
of sixes on the red dice, and $Y$ the number of twos on all three dice. Find the
joint distribution and $\Pbb(X>0\mid Y=1)$. Are $X$ and $Y$ independent?
Can you answer the last question without calculating the entire joint distribution?
\end{problembook}
\begin{problembook}
A fair coin is tossed independently four times. Let $X$ be the number of heads
in the first three tosses, and $Y$ the number in the last two.
\begin{enumeratea}
  \item Find the joint and conditional distributions, $\Pbb(X=Y)$ and $\Pbb(X>1\mid Y=1)$.
  \item Find $\Ebb(X\mid Y)$ and $\Ebb(Y\mid X)$ and their distributions.
  Verify that $\Ebb[\Ebb(X\mid Y)]=\Ebb X$.
  \item Find $\Cov(X,Y)$ in two ways: from the joint distribution, and by writing
  $X$ and $Y$ as sums of indicators.
\end{enumeratea}
\end{problembook}
\begin{problembook}
Let $X$ be uniformly distributed on $\{-1,0,1\}$ and $Y=X^2$.
\begin{enumeratea}
  \item Show that $\Cov(X,Y)=0$, but $X$ and $Y$ are dependent.
  \item Now let $X$ and $Y$ take values in $\{0,1\}$. Prove that $\Cov(X,Y)=0$
  is equivalent to independence. Why does this not contradict the preceding example?
  \item For random variables $X,Y$ with finite positive variances, prove that
  $|\Cor(X,Y)|\le1$. When does equality hold?
\end{enumeratea}
\end{problembook}
\begin{problempencil}
One box contains 4 white and 6 black balls, and another contains 7 white and
3 black balls. From each box we draw one ball uniformly, independently of the
other draw. We add one white ball to the two drawn balls and place all three in a third box.
\begin{enumeratea}
  \item What is the probability that a uniformly chosen ball from the third box is white?
  \item Draw from the third box with replacement until the first white ball.
  For a fixed box composition, draws are independent and uniform. Let $X$ be the
  number of draws. Find $\Ebb X$. Does $X$ have a geometric distribution?
  \item Find $\Pbb(X>k)$ for every integer $k\ge0$, and prove that
  $\Pbb(X>k+1\mid X>k)$ increases with $k$. What do we learn about the box's
  composition after a sequence of black balls?
\end{enumeratea}
\end{problempencil}
\subsection*{Additional problems}
\begin{problembook}
Choose three numbers uniformly without replacement from $\{1,2,3,4,5\}$.
Let $X$ be the middle number in order of size, and $Y$ the smallest.
\begin{enumeratea}
  \item Find the joint and marginal distributions. Are $X,Y$ independent?
  \item Find $\Cov(X,Y)$ and $\Cor(X,Y)$.
  \item Find the distribution, expectation and variance of $X-2Y$.
\end{enumeratea}
\end{problembook}
\begin{problempencil}
Let $(X_1,\ldots,X_n)$ be a uniformly random permutation of $\{1,\ldots,n\}$,
where $n\ge2$, and let $S=X_1+\cdots+X_n$.
\begin{enumeratea}
  \item Find $\Ebb S$, $\Var S$, $\Ebb X_i$ and $\Var X_i$.
  \item Prove the variance-of-a-sum formula and use it to find $\Cov(X_i,X_j)$
  for $i\ne j$. Explain the sign of the covariance.
  \item For $1\le k\le n$, find $\Var(X_1+\cdots+X_k)$. Compare with the variance
  of $k$ independent draws with replacement.
\end{enumeratea}
\end{problempencil}
\begin{problempencil}
On each face of a blank die, independently write a uniformly chosen number from
1 to 6. Then roll the resulting die $n\ge2$ times. The faces rolled are independent,
uniformly selected and independent of the labels. Let $T_k$ be the number seen on roll $k$.
\begin{enumeratea}
  \item What is the distribution of $T_1$? Find $\Cor(T_1,T_2)$. Are $T_1,T_2$ independent?
  \item If $X=T_1+\cdots+T_{n-1}$ and $Y=T_2+\cdots+T_n$, find $\Cor(X,Y)$.
  \item How do the two correlations change if the labels are independent and uniform
  on $\{199,\ldots,999\}$? What if the labels are independent with an arbitrary
  common distribution having finite positive variance?
\end{enumeratea}
\end{problempencil}
\begin{problem}[P\'olya's urn]
An urn contains one red and one black ball. At each step, choose a ball uniformly,
replace it and add another ball of the same colour. Let $X_n$ indicate a red ball
on draw $n$, and let $W_n=X_1+\cdots+X_n$. Find $\Pbb(X_n=1)$, $\Cov(X_i,X_j)$
for $i\ne j$, and the distribution of $W_4$. What distribution do you conjecture for $W_n$?
\end{problem}
\subsection*{Challenges}
\begin{problemstar}[Pairwise independence]
Let $q$ be prime, and let $A,B$ be independent and uniform on $\{0,\ldots,q-1\}$.
All subsequent calculations are modulo $q$. For $t\in\{0,\ldots,q-1\}$, define
\[ X_t=A+tB\pmod q,\qquad X_\infty=B. \]
\begin{enumeratea}
  \item Prove that $X_0,\ldots,X_{q-1},X_\infty$ are uniform and pairwise
  independent, but not mutually independent. What do the values of two distinct coordinates determine?
  \item Let $N$ be the number of zero coordinates in this family.
  Find its distribution, expectation and variance.
  \item Calculate the variance using indicators as well. Is pairwise independence
  enough for variances to add?
\end{enumeratea}
\end{problemstar}
\begin{problemstar}
For random variables $X,Y$ with finite second moments and a random variable $Z$, prove that
\[
\Cov(X,Y)=\Ebb\!\left[\Cov(X,Y\mid Z)\right]
+\Cov\!\left(\Ebb[X\mid Z],\Ebb[Y\mid Z]\right).
\]
What formula results when $Y=X$? Now let $Z\sim\Ber(1/2)$, let $U$ be uniform
on $\{-1,1\}$, let $Z,U$ be independent, and define
\[ X=3Z+U,\qquad Y=3Z-U. \]
\begin{enumeratea}
  \item Find the conditional covariance and correlation for fixed $Z$.
  \item Find $\Cov(X,Y)$ and $\Cor(X,Y)$.
  \item How do you explain the difference in signs? Relate your answer to the two terms in the formula.
\end{enumeratea}
\end{problemstar}
\subsection*{Part II. Generating functions and random sums}
\begin{problempencil}[Test 1, CS 2024]
For a nonnegative integer-valued random variable $X$, write $g_X(s)=\Ebb[s^X]$, $0\le s\le1$.
\begin{enumeratea}
  \item Let $X$ be uniform on $\{1,\ldots,6\}$. Find $g_X$ and express
  $\Ebb X$ and $\Var X$ in terms of its derivatives.
  \item If $X,Y$ are independent nonnegative integer-valued random variables,
  prove that $g_{X+Y}=g_Xg_Y$.
  \item[$(\mathrm{c})^\star$] Is the converse true?
\end{enumeratea}
\end{problempencil}
\begin{problembook}
Let $X,Y$ be independent. Use generating functions to find the distribution of
$X+Y$ when $X\sim\Bin(n,p)$ and $Y\sim\Bin(m,p)$, and then when
$X\sim\Poi(\lambda)$ and $Y\sim\Poi(\mu)$. Similarly find the distribution of
$G_1+\cdots+G_r$, where $G_1,\ldots,G_r$ are independent and $G_i-1\sim\Geom(p)$,
so each $G_i$ counts the trials up to the first success. Why does equality of
generating functions determine the distribution uniquely? Can the sum of independent
$\Bin(n,p)$ and $\Bin(m,q)$ variables be binomial when $n,m\ge1$, $0<p,q<1$ and $p\ne q$?
\end{problembook}
\begin{problempencil}
The number of visitors to a stadium is $N\sim\Poi(\lambda)$. Each visitor selects
one of ten entrances uniformly, independently of everyone else and of $N$.
\begin{enumeratea}
  \item Find the distribution, expectation and variance of the number using the first entrance.
  \item More generally, let $N,X_1,X_2,\ldots$ be independent nonnegative
  integer-valued random variables, with $(X_i)$ identically distributed.
  For $S=X_1+\cdots+X_N$, setting $S=0$ when $N=0$, prove that
  \[ g_S(s)=g_N\bigl(g_{X_1}(s)\bigr). \]
  If $N,X_1$ have finite second moments, find $\Ebb S$ and $\Var S$.
  \item Apply the result to the stadium problem. Also obtain the distribution
  in (a) using the law of total probability.
\end{enumeratea}
\end{problempencil}
\begin{problembook}
$A$ and $B$ play independent games. In each game, $A$ wins with probability $2/3$
and $B$ with probability $1/3$. Play stops when someone wins two consecutive games.
Let $T$ be the number of games played. Find $\Ebb T$ and the probability that
$A$ wins the contest. What information about the past suffices to describe future
play? Also find the generating function of $T$.
\end{problembook}
\subsection*{Additional problems}
\begin{problempencil}
Two teams each take one penalty in each round. All attempts are independent,
with scoring probabilities $0.75$ for the first team and $0.80$ for the second.
All five regular rounds are played. If tied afterwards, play continues until,
in one round, one team scores and the other does not.
\begin{enumeratea}
  \item What is the probability that in each of the first five rounds both teams
  have the same outcome for their penalty?
  \item What is the probability that play continues after the fifth round? After the tenth?
  \item Find the expected total number of penalties taken.
\end{enumeratea}
\end{problempencil}
\begin{problem}
Lottery tickets are numbered from 000000 to 999999. Find the probability that a uniformly chosen ticket:
\begin{enumeratea}
  \item has digit sum 21;
  \item has equal sums of its first and last three digits;
  \item has a sum of its first three digits that exceeds that of its last three by 2.
\end{enumeratea}
It suffices to express the answers as coefficients of a polynomial.
\end{problem}
\begin{problempencil}
We telephone $N\sim\Poi(\lambda)$ people with an offer. Person $i$ would accept
on the $K_i$th call, where $K_1,K_2,\ldots$ are independent, identically distributed
and independent of $N$, with $\Pbb(K_i=n)=f(n)$ for $n\ge1$ and
$\sum_{n\ge1}f(n)=1$. Let $X_n$ be the number of people who would accept on call $n$.
\begin{enumeratea}
  \item Prove that $(X_n)_{n\ge1}$ are independent and $X_n\sim\Poi(\lambda f(n))$.
  \item We stop at the first round with no new acceptance. If
  $T=\inf\{n\ge1:X_n=0\}$ and $S=X_1+\cdots+X_T$, prove that
  \[ \Ebb S=\lambda\Ebb[F(T)],\qquad F(k)=\sum_{j=1}^k f(j). \]
  Why can we not directly apply the expectation formula for a random sum whose
  number of summands is independent of the summands?
\end{enumeratea}
\end{problempencil}
\begin{problempencil}
Two standard dice are rolled independently. Can we label their faces with other
positive integers so that independent rolls have the same sum distribution?
Repeated labels are allowed. Find a construction different from the two standard dice.
\end{problempencil}
\subsection*{Challenges}
\begin{problemstar}\label{prob:week4-penney}
We toss a fair coin independently, recording outcomes as $H,T$. For a word
$w=w_1w_2w_3$, let $\tau_w$ be the toss number when it first appears as three
consecutive outcomes. Consider $a=HHT$ and $b=THH$.
\begin{enumeratea}
  \item Choose states retaining enough information about the appearance of the
  two words, and draw the possible transitions.
  \item Find $\Pbb(\tau_a<\tau_b)$ and $\Ebb(\tau_a\wedge\tau_b)$.
  \item Prove that $\tau_a\wedge\tau_b<\infty$ almost surely, without using the calculated expectation.
\end{enumeratea}
\end{problemstar}
\begin{problemstar}[Branching process]
A population starts with one individual. Each individual has 0 children with
probability $1-p$ and 2 children with probability $p$, where $0<p<1$, independently
of all other individuals in all generations. Let $Z_n$ be the number of individuals
in generation $n$, with $Z_0=1$, and let
\[ f(s)=1-p+ps^2,\qquad T=\inf\{n\ge0:Z_n=0\}. \]
\begin{enumeratea}
  \item Prove that $g_{Z_n}=f^{\circ n}$.
  \item Show that $q_n=\Pbb(T\le n)=f^{\circ n}(0)$ and that $(q_n)$ is increasing.
  \item Prove that $q=\Pbb(T<\infty)$ is the smallest solution of $f(q)=q$ in
  $[0,1]$, and find it as a function of $p$.
  \item Find $\Ebb Z_n$. When $p=1/2$, it is constant, yet extinction occurs almost surely.
  How do you explain this?
\end{enumeratea}
\end{problemstar}
\subsection*{Independent research}
\begin{researchtask}[Penney's game]
Two players each choose a different word from $\{H,T\}^3$. We then toss a fair
coin independently until one of the words appears. The player whose word appears first wins.
\begin{enumeratea}
  \item Generalise the model in Problem~\ref{prob:week4-penney} and write a program
  that calculates exactly, for each of the 56 ordered pairs $(a,b)$,
  \[ p(a,b)=\Pbb(\tau_a<\tau_b),\qquad e(a,b)=\Ebb(\tau_a\wedge\tau_b). \]
  Use rational arithmetic and check at least one pair by hand.
  \item For each first word $a$, find the best responses $b$ for the second player.
  Does the player choosing second have an advantage? Formulate a general rule
  for the best response and prove it.
  \item Draw a directed graph with edge $a\to b$ when $b$ is a best response to
  $a$. Find a cycle and explain what it shows about comparisons between words.
\end{enumeratea}
\end{researchtask}
\noindent\textbf{Lean.} Formalise in Lean 4 with Mathlib and prove that, for
finite sets $S,T$ and functions $a:S\to\mathbb R$ and $b:T\to\mathbb R$,
\[
\sum_{s\in S}\sum_{t\in T}a(s)b(t)
=\left(\sum_{s\in S}a(s)\right)\left(\sum_{t\in T}b(t)\right).
\]
Explain how this identity is used in proving $g_{X+Y}=g_Xg_Y$ for independent
nonnegative integer-valued random variables $X,Y$ with finite ranges.

\end{document}
